\chapter{Lehre}

Das Projekt ist auch ein Lerngegenstand. Dieses Kapitel beschreibt, wie es in Kursen eingesetzt werden kann.

\section{Kurse}

\startitemize
\item \textbf{DUA25, Datenstrukturen und Algorithmen:} Hashing, Merkle-Graphen, Myers-Diff und Graphsuche.
\item \textbf{AISE, AI Assisted Software Engineering:} Agenten-Workflows, Branch-Regeln, Freigaben und Risikostufen.
\stopitemize

\section{Übungen je Phase}

Jede Phase lässt sich als Übung formulieren. Studierende implementieren einen Teil selbst und vergleichen ihn mit einer Referenzimplementierung. Die Phase 4 eignet sich besonders, weil Freigaben und Selbstfreigabeverbot eine Prüfung durch Agenten und Menschen verlangen.

\section{Offene Fragen}

\startitemize
\item Lizenz des Lehrmaterials und Umgang mit fremden Quellen.
\item Bewertungskriterien für Code, der mit Agenten entstanden ist.
\item Trennung von Studierendennamespaces und Dozierendenrolle.
\stopitemize
